% Proposed replacement for Section 8 of main.tex (lines 618--799 of the fourth draft as of 03:31 UTC, 24 November 2026;
% that text is identical to lines 617--798 of the 03:25 version).
% Same numbers, same failed checks, same caveats; reordered. Labels kept: s:evid, s:local, tab:local,
% s:time, s:cross. New labels: tab:runs, tab:cross.
\section{Measurements}\label{s:evid}
This section reports Monte Carlo estimates for $\SU(2)$. Nothing in it is used in a proof. Each subsection
asks a question, says in a sentence what was found, and then gives the setup, the numbers, the checks that
failed and what the numbers do not show. The design of the main runs is collected in Table~\ref{tab:runs}.

\emph{What a measurement can show.} Each measurement tries particular test functions in an inequality that an
open hypothesis, or the auxiliary statement, would need for all functions. The exact value for a trial
function limits a constant from one side. What is reported is an estimate of that value with a statistical
error, which is not a bound, and no measurement can show that a hypothesis holds.

\emph{Conventions.}
\begin{itemize}
\item A \emph{stream} is an independent run of the sampler with its own seed, started hot or cold in
alternation; the stream is the unit of replication.
\item Errors are delete-one-stream jackknife standard errors, a digit in parentheses is the standard error in
the last digit, and the number of streams behind each error is stated with it.
\item Where a test function was selected from a span, the selection used one set of streams and the reported
value comes from other streams; the error is then that of the selected function with its coefficients fixed,
and does not include the variability of the selection.
\end{itemize}

\emph{Limits common to all of them.}
\begin{itemize}
\item One coupling per lattice; lattices of side $5$ to $10$ with periodic time; heat-bath sampling.
\item Thermalisation was checked only by comparing hot with cold starts and first with second halves, and
some of these checks failed; they are named where they bear on a number.
\item Couplings were taken from a calibration in~\cite{RepoData} and not recalibrated, so every statement at
``matched'' couplings inherits it.
\item Programs, seeds, the rules fixed before each run, the rules that failed and summary files are
in~\cite{RepoData}. Per-stream data are deposited for Section~\ref{s:cross} and for the single-mark runs; for
the others they are regenerated from the seeds, which reproduces them statistically but, across rebuilds of
the sampler, not byte for byte.
\end{itemize}

\begin{table}[t]
\centering\small
\begin{tabular}{@{}>{\raggedright\arraybackslash}p{0.19\textwidth}>{\raggedright\arraybackslash}p{0.27\textwidth}>{\raggedright\arraybackslash}p{0.24\textwidth}>{\raggedright\arraybackslash}p{0.22\textwidth}@{}}
\toprule
Runs & Lattice and coupling & Streams and length & Use of the streams \\
\midrule
Frozen backgrounds A and B (Section~\ref{s:local}) &
side $N=6$, $T=12$, $\beta=3.276133577$, a point on a line of constant renormalised coupling &
each background has one run of $32$ streams ($576$ sweeps discarded, then $1152$ measurements six sweeps apart) &
selection on streams $0$ to $15$; reported value on streams $16$ to $31$ \\
\addlinespace
Carry-over in time (Section~\ref{s:time}) &
same lattice, coupling and frozen slice (background A); the later slices are not frozen &
$32$ fresh streams &
$16$ for fitting and $16$ held out \\
\addlinespace
First span (Section~\ref{s:cross}) &
sides $N=5$ and $N=10$ at couplings $3.196581051$ and $3.493705679$ (matched at a fixed renormalised
coupling); time extents two and four times the side; a common physical lag (one and two time steps) &
$16$ streams per lattice, $500$ measurements each &
half of the streams used for selection; reported value on the $8+8$ held-out streams \\
\addlinespace
Second study (Section~\ref{s:cross}), designed after the first was seen &
as for the first span &
new streams ($16$ per lattice, $2000$ measurements each) &
tests fixed on the selection streams; reported values on the $8+8$ held-out streams \\
\addlinespace
Longer ensembles (Section~\ref{s:cross}) &
ensembles of twice the time extent &
$16+16$ streams; $1000$ measurements per stream on side $10$ &
no selection; they share no stream with the selection \\
\bottomrule
\end{tabular}
\caption{Design of the main runs. The runs on displaced backgrounds and the single-mark runs of
Section~\ref{s:local} are described there and in~\cite{RepoData}.}\label{tab:runs}
\end{table}

\subsection{The local quotient on a frozen coarse background}\label{s:local}
\emph{Question.} On a coarse background with holonomy $-I$, how small can the local quotient
$\mathcal J/\mathcal V$ be?

\emph{What we found.} On the backgrounds tried, the paired differences show the selected quotient still
decreasing at the largest span of trial functions, although the rule fixed before the run called the decrease
stopped on A; no lowest value is indicated. For these spans the quotient is smaller on the background with
holonomy $-I$ than on the other. All values are estimates for trial functions, not bounds.

\emph{Setup.} Lattice, coupling and streams are in Table~\ref{tab:runs}. The coarse links of one slice are
held fixed (a \emph{frozen background}) and the fine links are sampled around them.
\begin{itemize}
\item \emph{Background~A.} All coarse links are the identity except two. One is $-I$; of the four coarse
spatial plaquettes through it, three have holonomy $-I$ and the fourth also contains the second altered link.
That link, adjacent to the first, is the quaternion $\mathbf i$, which supplies a reference direction; four
coarse plaquettes, the fourth among them, then have holonomy $\pm\mathbf i$.
\item \emph{Background~B.} A single configuration: the first of $60$ drawn from the unconditioned measure
whose coarse plaquettes are all farther than $\pi/4$ from $-I$ ($40$ qualified), a rule fixed in advance.
\item \emph{Marks.} A \emph{mark} is a gauge-invariant function of the fine plaquettes inside one coarse
plaquette; the list of marks is in~\cite{RepoData}.
\item \emph{Spans.} Four nested spans were fixed before the runs: the $15$ marks on the three $-I$
plaquettes; $49$ marks on the seven coarse plaquettes that contain an altered link; $181$ and $308$ marks
that add the $33$ plaquettes sharing a coarse edge with those.
\item \emph{Selection.} For each span the linear combination with the smallest quotient is selected on the
selection streams, and Table~\ref{tab:local} reports its quotient on the held-out streams. Smaller means
slower relaxation under the local updates.
\end{itemize}

\begin{table}[ht]
\centering\small
\begin{tabular}{@{}llcccc@{}}
\toprule
 & & \multicolumn{4}{c}{Number of marks in the span} \\
\cmidrule(l){3-6}
Background & & 15 & 49 & 181 & 308 \\
\midrule
A & held-out quotient  & .0208(5) & .0167(4) & .0159(4) & .0154(4) \\
  & quotient on selection streams & .0213 & .0172 & .0163 & .0157 \\
  & decrease from the span to its left &  & 20\% & 5\% & 3\% \\
\addlinespace
B & held-out quotient  & .0589(8) & .0533(5) & .0510(4) & .0428(4) \\
  & quotient on selection streams & .0587 & .0514 & .0482 & .0410 \\
  & decrease from the span to its left &  & 10\% & 4\% & 16\% \\
\bottomrule
\end{tabular}
\caption{Local quotient $\mathcal J/\mathcal V$ (dimensionless) of the combination of marks selected on $16$
streams and evaluated on the $16$ other streams of the same run. A: frozen background with three coarse
plaquettes of holonomy $-I$. B: one frozen background drawn from the unconditioned measure, every coarse
plaquette at round distance greater than $1.1$ from $-I$. Errors are over the $16$ held-out streams with the
selected coefficients fixed. The entries of a row are computed on the same streams and are strongly
correlated; the two rows come from independent runs. All entries are estimates, not bounds.}
\label{tab:local}
\end{table}

\emph{Unselected functions, for comparison.} In the same runs and on all $32$ streams, the unselected average
of the four plaquette marks of one coarse plaquette gives $.0473(5)$ on A and $.0893(9)$ on B, and a single
plaquette mark gives $.0611(4)$ on A.

\emph{{\upshape(i)} The selected combination is collective.} In the $308$-mark span on A, $81$ per cent of
the sum of its squared coefficients, each scaled by the standard deviation of its mark, belongs to the three
coarse plaquettes with holonomy $-I$.

\emph{Caveats.} Four exact linear relations among these marks were removed and the coefficients of least norm
taken, so the share depends on that convention. The marks are also correlated, so this is a statement about
coefficients and not a decomposition of the variance. On B the selected direction is less stable: for the
$49$-mark span the directions selected on the two halves of the streams have cosine $0.41$.

\emph{{\upshape(ii)} No lowest value is indicated.} The spans are nested, so the smallest quotient of a span
cannot increase as it grows; what is measured is the size of the steps.
\begin{itemize}
\item \emph{By the rule fixed before the run.} The rule compared successive entries as if independent and
called the decrease stopped when they differ by less than two combined standard errors. By that rule the
decrease stopped after $49$ marks on A and did not stop on B.
\item \emph{By paired differences, which the rule did not use.} The entries of a row share their streams,
however, and the paired differences are $38$, $15$ and $12$ standard errors on A and $9$, $6$ and $15$ on B.
On B the last step is the largest.
\end{itemize}
The ratio of the two rows, about $2.8$, is for these spans only.

\emph{{\upshape(iii)} Displaced backgrounds: a few points, not an interval.} Here the frozen holonomy is at a
stated round distance from $-I$.
\begin{itemize}
\item \emph{The same $15$ marks.} At distance $1/\beta$ they give $.0207(2)$, against $.0216(5)$ at $-I$ in a
separate, earlier run of the same design.
\item \emph{The single plaquette mark.} The ratio of the quotient at distance $2/\beta$ to that at $-I$ is
$1.016(11)$ in a first run and $1.004(11)$ in a second run whose rules were registered beforehand; at
$4/\beta$ it is $1.075(11)$.
\item \emph{The four-mark average.} The ratio at $2/\beta$ is $1.034(13)$.
\end{itemize}

\emph{Checks that failed.} Two checks failed:
\begin{itemize}
\item the hot-against-cold check of the first run at $2/\beta$ ($2.7$ standard errors; it had been reported
as passed because of an error in the reduction program, since repaired);
\item one of the four burn-in checks of the second run, in which all four drifts have the same sign.
\end{itemize}

\emph{Reading.} These points indicate that quotients like those at $-I$ are also found at distances of order
$1/\beta$; they do not determine how far.

\emph{Comparison with the proved bound.} The bound of Theorem~\ref{t:ret} at this coupling is
$p_{6\beta}(\theta)\approx0.02$, $0.15$, $0.39$ and $0.93$ at $\theta=0.1$, $0.2$, $1/\beta$ and $2/\beta$: it
is small only well inside the distances tried.

\subsection{Carry-over in time of the locally slow combination}\label{s:time}
\emph{Question.} Is the combination that the local updates relax most slowly also slow to change in time?

\emph{What we found.} A clause of the rule fixed before the runs is met formally, but we draw no verdict
from it. For
this one combination on this one frozen slice, the linear predictors that were tried recover at most a few
per cent of its variance one slice later. This is not a measurement of Hypothesis~\ref{h:mix}.

\emph{Setup.} Lattice, background and streams are in Table~\ref{tab:runs}. The $308$-mark combination
selected above was regressed on the same $308$ marks one time slice later, in two gauge-invariant forms:
\begin{itemize}
\item the \emph{later-slice form}, with reference directions taken on the later slice;
\item the \emph{transported form}, with the later marks transported back along time-like links to the frozen
reference (a mark of this form is not a function of the later slice alone).
\end{itemize}
The rule fixed before the runs had three outcomes:
\begin{itemize}
\item persistence, if the held-out coefficient of determination exceeds $0.1$ and exceeds that on
background B by three standard errors;
\item no persistence, if it is below $0.02$ with error below $0.01$;
\item otherwise unresolved.
\end{itemize}

\emph{Result under the rule.} The values are $-0.013(2)$ and $-0.002(2)$ in the two forms (errors over the
held-out streams with the fitted coefficients fixed; the uncertainty of the fit is not included).

\emph{Why no verdict is drawn.} With $308$ regressors and $18432$ fitting measurements an uncorrelated target
gives about $-0.017$, and the same fit at a lag of three slices, where no correlation is seen, gives
$-0.016(2)$ and $-0.019(2)$. The second clause of the rule is therefore met formally, but with this offset it
cannot exclude a coefficient of determination up to about $0.03$, and we draw no verdict from it.

\emph{Outside the rule} (chosen after the data were seen).
\begin{itemize}
\item A fit on $15$ marks in the transported form gives $0.011(2)$ at lag $+1$ and $0.005(2)$ at lag $-1$,
which estimate the same quantity.
\item The correlation of the combination with its own copy one slice later (all $32$ streams) is $0.083(6)$
in the transported form and $0.022(4)$ in the later-slice form, and $0.016(6)$ in the transported form at two
slices.
\item On background B the combination selected there gives $0.032(5)$ and $-0.003(6)$.
\item The unselected four-mark average gives $0.076(5)$ on A and $0.009(6)$ on B in the transported form, so
the larger correlation on A is not particular to the selected combination.
\end{itemize}
The result files of this run are not in the deposit; these numbers are those of the run's
report~\cite{RepoData}.

\emph{What this does not show.}
\begin{itemize}
\item The later slice is not frozen, so its coarse background is not the one for which the combination was
selected.
\item The test is a correlation conditional on the coarse links of one slice, not a matrix element of the
transfer operator.
\item It does not bound the best prediction from the later slice, and it does not exclude other combinations
that propagate slowly.
\item The local quotient and a correlation over one time step are not on a common scale, so the two are not
compared.
\item The distinction between Hypothesis~\ref{h:mix} and the auxiliary statement is the logical one of
Section~\ref{s:arch}; this measurement neither supports nor contradicts either.
\end{itemize}

\subsection{Comparison across a factor of two in the mesh}\label{s:cross}
\emph{Question.} Does a function of the coarse links decay in time on the fine lattice as it does on the
coarse one?

\emph{What we found.} Without normalisation, the comparison of energies fails for the selected test as a
matter of normalisation. With normalisation, a test of the second study has a quotient ratio below
one, on held-out streams and again on longer ensembles; the other tests fixed there do not, at the stated
margin. A discrepancy between programs is unexplained. None of this measures
Hypothesis~\ref{h:coarse}.

\emph{Setup.} Lattices, couplings and streams are in Table~\ref{tab:runs}. The tests are functions of the
coarse links of side $10$, and the same functions of the links of side $5$. For a test $f$ let
\[
\mathsf R(f)=1-C_f(\mathrm{lag})/C_f(0),\qquad \mathcal E_N(f)=C_f(0)-C_f(\mathrm{lag}),
\]
with $C_f$ its connected two-time correlation. So $1-\mathsf R$ is the correlation retained over the lag, and
$\mathcal E_N(f)$, taken on the lattice of side $N$, is its two-time energy.

\emph{{\upshape(i)} Unnormalised energies.} In a first span of $16$ slice-averaged functions, the test
selected to minimise the ratio of energies has, on the held-out streams,
$\mathcal E_{10}=0.00407(5)\,\mathcal E_5$. This is the ratio of its stationary variances, $0.0034$, times the
ratio of its quotients, $1.21$: the function has a much smaller variance on the finer lattice and decays
slightly faster there. An inequality $\mathcal E_{10}(f)\ge c\,\mathcal E_5(f)$ with $c$ of order one
therefore fails for this function on this pair as a matter of normalisation, and a comparison of deficits
must account for the change of variance. In this span no test had a quotient ratio below one.

\emph{{\upshape(ii)} Normalised quotients.} The second study used a span of $54$ functions ($48$ linearly
independent among the slice-averaged ones on the selection streams) and five classes of test. Of the $25$
tests fixed on the selection streams, one has a ratio below one on the held-out streams: the test selected to
minimise $\mathsf R_{10}/\mathsf R_5$ among slice-averaged functions. The other $24$ have ratio plus two
standard errors above one. The same fixed test was then evaluated on the longer ensembles.
Table~\ref{tab:cross} gives both evaluations.

\begin{table}[ht]
\centering\small
\begin{tabular}{@{}lcccc@{}}
\toprule
Evaluation of the one fixed test & $\mathsf R_5$ & $\mathsf R_{10}$ & ratio & ratio plus two standard errors \\
\midrule
held-out streams of the second study & $0.990(3)$ & $0.932(4)$ & $0.941(5)$ & $0.951$ \\
longer ensembles & $0.996(1)$ & $0.932(2)$ & $0.936(3)$ & $0.941$ \\
\bottomrule
\end{tabular}
\caption{The test selected to minimise $\mathsf R_{10}/\mathsf R_5$, evaluated with its coefficients fixed.
Estimates with jackknife errors over the streams of Table~\ref{tab:runs}; not bounds.}\label{tab:cross}
\end{table}

So this test retains about $7$ per cent of its correlation over the lag on the finer lattice and between
$0.4$ and $1$ per cent on the coarser. The rule fixed in advance reads off the last column on both pairs: for
this test, on this pair, at these couplings, the estimates do not support
$\mathsf R_{10}\ge\alpha\,\mathsf R_5$ with $\alpha>0.951$.

\emph{Caveat on that figure.} It is two jackknife standard errors over few streams, not a confidence limit at
a stated level, and it inherits the coupling match.

\emph{A check that failed: an unexplained discrepancy.} For a different statistic on the same finer lattice,
an earlier and shorter run by another program in~\cite{RepoData} gives $0.0041(7)$ where these ensembles give
$0.0016(1)$, $3.4$ combined standard errors apart. The cause is not determined.

\emph{What this does not show.}
\begin{itemize}
\item Hypothesis~\ref{h:coarse} compares the smallest deficits on the two lattices, not the values of one
test, and $\varepsilon_i$ is an allowed loss in that comparison. One test cannot measure it, and this result
is neither evidence for nor against the hypothesis.
\item The lattices have periodic time of finite extent and are not the vacuum; one pair of sides, one
renormalised coupling and one lag were used.
\item The first span contained no such test, so the outcome depends on the span, and a larger span could
change it again.
\end{itemize}
